% !TeX root = ../main.tex
% 本章按当前规格、CModel 与 Verilog 核验；英文稿按作者后续要求同步。
\section{面向配比探索的组合式点积微架构}
\label{sec:architecture}
\label{sec:brick_definition}

DP-Puzzle 采用 Lane、Group 与 DP 三层组合式架构（图~\ref{fig:dp_architecture}）。Lane 共享元素乘法与局部归约，Group 融合缩放因子并输出统一表示的部分和，DP 完成全局对齐、累加与输出处理。设计的核心是以统一的 Group 接口连接多精度共享通路与低精度专用通路：保留完整格式支持，通过增配专用 Group 调整各精度的峰值算力，同时共享末端累加、规格化与舍入。九种配置的硬件资源在设计阶段确定，运行时按请求切换计算精度。

\subsection{符号与记号约定}
\label{sec:notation_conventions}

本文以 T、H、Q、E 表示四类精度：T 为 TF32，H 为 FP16/BF16，Q 为 MXFP8（E4M3/E5M2）与 MXFP6（E2M3/E3M2），E 为 MXFP4、NVFP4 与 HiF4。后文用这些简称描述各类精度的计算通路与算力配比。

定点数采用 $\mathrm{Q}I.F$ 与 $\mathrm{UQ}I.F$ 记号：$I$ 为不含符号位的整数位数，$F$ 为小数位数~\cite{ti2003qformat}。$\mathrm{Q}I.F$ 表示二进制补码有符号数，总位宽为 $1+I+F$；$\mathrm{UQ}I.F$ 表示无符号数，总位宽为 $I+F$。例如，UQ2.26 为含 2 位整数与 26 位小数的 28 位无符号数。

计算并行度记为 $\mathrm{T}x_T\mathrm{H}x_H\mathrm{Q}x_Q\mathrm{E}x_E$，各下标变量表示对应模式每周期处理的 A/B 元素对数，每对产生一个乘积。例如，T1H2Q4E8 表示 T/H/Q/E 模式每周期分别计算 1/2/4/8 个乘积。各精度模式互斥，上述数值不能相加作为同时可用的吞吐。一个 Lane 支持 T1H2Q4E8，两个 Lane 构成一个 Group，四个 Group 构成 T8H16Q32E64 基线；每个 Lane 另支持四路 INT8 乘法，DP 支持 INT32 累加。各配置均采用八级流水，每拍接收一笔独立请求。

对于固定的 $m\times n$ DP 阵列，精度 $p$ 的原生点积长度为 $K_p$，时钟频率为 $f$。当每拍均有输入且阵列充分利用时，峰值算力为
\begin{equation}
\mathrm{Peak}_p=2mnK_pf\quad(\mathrm{FLOP/s}).
\label{eq:array_peak_throughput}
\end{equation}
一次 MAC 计为两次浮点运算，因此 $K_H:K_Q:K_E$ 直接决定同构阵列的峰值配比。后续设计通过改变低精度模式的 $K_p$ 构造候选，并分别表征其面积、功耗与供数开销。

\begin{figure}[t]
\centering
\includegraphics[width=\columnwidth]{baseline_arch_singlepage.pdf}
\caption{缩放粒度驱动的 Lane--Group--DP 划分。HiF4 每 8 个元素共享一个二级微指数，对应一个 Lane 的 8 路 E 类乘积；NVFP4 每 16 个元素共享一个尺度，对应两个 Lane 组成的 Group。更大的量化块向多个 Group 广播同一尺度，保持原格式的缩放粒度。}
\label{fig:dp_architecture}
\end{figure}

\subsection{Lane：跨精度共享与低精度定点化}
\label{sec:unpack_recode}
\label{sec:full_brick}

\textbf{统一解包与乘法复用。}图~\ref{fig:data_unpack} 给出各格式的内部表示。T/H 保留指数与有效数，Q 映射为内部 E5M3 表示；该映射通过指数调整和小数位补零保持原值。MXFP4/NVFP4 将元素指数融入 UQ3.1 幅值，HiF4 将一级微指数并入 UQ2.2 幅值。Lane 采用两组可复用的部分积单元：TF32/FP16 按位权组合部分积，Q/E 则将乘法叶子用于独立的低精度乘积，避免为每种格式复制完整乘法通路。

\begin{figure}[t]
\centering
% 图稿待作者统一数量标签：按完整 Lane 标注时，FP16 为 2x，E 为 8x；
% 右侧 10/16/4 位描述单个字段组，每个 Lane 包含两个字段组。
\includegraphics[width=\columnwidth]{unpack_layout.png}
\caption{各格式到共享计算字段的映射。Q 统一指数与有效数表示，E 将元素指数或一级微指数融入定点幅值；符号与特殊值状态独立传递。}
\label{fig:data_unpack}
\end{figure}

\textbf{按精度组织对齐与归约。}T/H/Q 使用组内共同的最大乘积指数，通过四路共享右移器完成局部对齐，幅值窗口为 UQ2.26。E 类在解包后采用固定坐标，无需逐元素指数比较与宽移位：八个乘积先两两合并，再接入共享的四路归约树。HiF4 的二级微指数由 Lane 部分和之后的小范围移位消化，避免引入跨 Lane 的 E 类指数对齐。由此，低精度计算可以复用乘法与加法资源，同时旁路高精度所需的对齐通路。

\subsection{Group：规约后尺度融合与统一部分和}

\textbf{将尺度运算移至组内规约之后。}两个 Lane 共享同一组尺度，Group 先合并其部分和，再对幅值施加两侧尺度有效数的乘积。NVFP4/HiF4 共用一套 $14\times8$ 无符号乘法器，处理 14 位组幅值与 8 位尺度有效数乘积；MX 的幂次尺度仅计入指数。尺度乘法由此按 Group 执行，无需在每个元素乘积上重复计算。物理 Group 的划分不改变量化块大小，同一量化块的尺度按需广播到多个 Group。

\textbf{分离数值坐标与幅值处理。}各类 Group 统一输出符号、31 位幅值及指数信息。令幅值为 $\mathrm{SoP}$、位宽为 $M$、符号为 $s$，原始坐标 $\mathrm{Anchor}$ 表示幅值最高物理位的指数，则有限组结果的数值为
\begin{equation}
v=(-1)^s\mathrm{SoP}\,2^{\mathrm{Anchor}-(M-1)}.
\label{eq:anchor_recompute}
\end{equation}
Anchor 由局部乘积指数、尺度指数及格式相关的定点布局修正提前计算，与幅值归约并行。对有限非零部分和，前导零检测仅产生比较指数 $e=\mathrm{Anchor}-\operatorname{LZC}_M(\mathrm{SoP})$，不在 Group 内移动幅值。该接口既保留原始数值坐标，又为后端提供实际最高有效位的指数，使不同专用 Group 能共用同一对齐与累加通路。

\subsection{DP：一次净移位与共享融合累加}

累加项 $C$ 在组内规约和尺度融合之后才加入全局指数比较与累加，避免其指数提前影响 E 类的局部定点计算。Group 使用经前导零校正的指数 $e$；$C$ 使用解码后的有效指数，保留 FP32 非正规数的指数约定，不做前导零校正。有限非零参与源共同确定全局参考 $E_{\max}$，零值与特殊值不参与数值 MAX。

每路原幅值低位补零至 39 位窗口后，按净移位量 $\Delta=\mathrm{Anchor}-E_{\max}$ 执行一次双向移位：正值左移，负值右移并向零截断。该操作将 Group 的前导零规格化与全局对齐合并，无需先移动幅值再重新对齐。恢复符号后，各路进入 carry-save 压缩网络，最后通过一次进位传播加法得到总和。中间部分和不转为 FP32，仅在最终输出时规格化并采用最近偶数舍入。INT8 复用乘法资源，与 INT32 累加项共用末端加法通路，旁路浮点对齐和输出规格化舍入。

\subsection{低精度专用 Group}
\label{sec:specialized_midgroups}

专用化保留统一的 Group 输出接口，并按目标精度裁剪解包、乘法、对齐和寄存器，而非只改变输入格式选择。本文提供三类可组合单元。

\textbf{E 专用单元。}E8 Lane 仅保留低精度解包、八个 $4\times4$ 乘法器及窄归约树，HiF4 的局部微指数在 Lane 内完成融合。两个 Lane 构成 E16 Group，复用组级尺度乘法，补零后接入公共部分和接口。该路径无需 T/H/Q/INT8 的解码、元素指数比较与宽对齐逻辑。

\textbf{Q/E 共享单元。}Q4E8 Lane 以八个 $4\times4$ 乘法器支持两类精度：Q 模式启用四个，E 模式启用全部八个。两个 Lane 构成 Q8E16 Group；Q 保留指数比较与宽对齐，E 旁路对齐并复用归约和尺度处理。因而 QE 增配仍需承担 Q 的宽对齐与规约成本，不能按 E 专用单元的成本线性放大。

\textbf{Q 专用单元。}Q4 Lane 与 Q8 Group 仅保留 Q 的乘法、指数对齐及归约，删除 E 相关的配对、微指数与尺度有效数乘法，用于单独提高 Q 吞吐。三类专用单元分别支持 E 定向增配、Q/E 共同增配与 Q 定向增配。

\subsection{DP 内算力配比配置}
\label{sec:expansion}

四个多精度 Group 构成 T8H16Q32E64 基线。在此基础上增加 2、4、8、12 个 E16 Group，分别得到 $1:2:6$、$1:2:8$、$1:2:12$、$1:2:16$；增加 4 个 Q8 Group 得到 $1:4:4$；增加 4、8、12 个 Q8E16 Group，分别得到 $1:4:8$、$1:6:12$、$1:8:16$。连同 $1:2:4$ 基线，共形成九种全精度配比，均保留 TF32、FP16/BF16 与 INT8 通路。

新增 Group 计算同一输出的额外 A/B 片段，与基础 Group 使用统一的全局参考指数，在 carry-save 网络内合并部分和。全局 MAX、各源对齐移位器和规约网络随组数扩展；$C$ 只加入一次，最终进位传播加法器、规格化与舍入仍保持单份。网络末端按请求精度选择基础或扩展后的两行，再送入唯一的进位传播加法器，避免在分支合并之前引入串联的进位传播。累加位宽随源数增长，保留对齐后各项求和所需的增长位。

各配置均保持八级流水与 II=1。扩展组按对应请求的精度更新寄存器，非参与组从指数比较、数值累加及特殊值归约中屏蔽。该组织使算力配比成为设计阶段的可组合参数，并将各候选的成本差异落实到专用前端、扩展对齐和规约资源，为后续面积、功耗与部署规模的联合分析提供基础。
